\section{Conclusion}\label{sec:conclusion}
THERMO-SMART est une plateforme de simulation thermique qui résout l'équation de la chaleur par éléments finis P1
avec FreeFem++, en régime stationnaire et transitoire, sur des parois multi-matériaux et sur une coupe de pièce en
2D, puis compare des solutions d'isolation, de protection solaire et de ventilation. Chaque étape a été validée avant
la suivante : ordre de convergence 2 en norme $L^2$ et 1 en norme $H^1$, températures exactes dans un mur de trois
couches, bilans de flux à la précision machine, contrôle analytique du cycle jour/nuit et critères de régime
périodique vérifiés sur chaque scénario.

Les principaux résultats, valables dans le cadre des hypothèses du modèle, sont les suivants :
\begin{itemize}
	\item les ponts thermiques aux jonctions et au plancher ajoutent environ 13~W/m au flux entrant de la pièce
	de référence (28~\% du total), presque indépendamment de l'isolant des murs ;
	\item supprimer l'isolant double les degrés-heures au-dessus de 30~\textdegree C ($+105$~\%), alors qu'un isolant
	plus performant les réduit de 24~\% ;
	\item une peinture claire sur le toit divise par deux l'excès de chaleur dû au soleil ($-50$~\%) ;
	\item la ventilation nocturne est la mesure la plus efficace : $-1{,}6$~\textdegree C en moyenne et $-90$~\% de
	degrés-heures.
\end{itemize}

Une interface web permet de modifier les paramètres et de relancer les calculs, et l'ensemble du dépôt (scripts,
configurations, définitions des cas tests, README) permet de reproduire chaque résultat. Les limites du modèle, en
particulier l'absence de fenêtres, d'apports internes et de comparaison avec des mesures, sont détaillées à la
section~\ref{sec:limites}, avec les pistes pour la suite : ventilation commandée, solutions contre les ponts
thermiques, données météorologiques réelles et, après validation, apprentissage automatique pour accélérer
l'exploration des configurations.

\section*{Références utiles}
\begin{itemize}
	\item FreeFem++ : \url{https://freefem.org/}
	\item Python : \url{https://docs.python.org/3/}
	\item NumPy : \url{https://numpy.org/doc/}
	\item Matplotlib : \url{https://matplotlib.org/stable/}
	\item Streamlit : \url{https://docs.streamlit.io/}
\end{itemize}